\documentclass[8pt]{extarticle}
\usepackage[utf8]{inputenc}
\usepackage[T1]{fontenc}
\usepackage[a4paper,landscape,left=0.3cm,right=0.3cm,top=1cm,bottom=0.3cm,headheight=12pt,headsep=4pt]{geometry}
\usepackage{hyperref}
\usepackage{fancyhdr}
\usepackage{extramarks}
\usepackage{ifthen}
\usepackage{listings}
\usepackage{xcolor}
\usepackage{tocloft}
\usepackage{multicol}
\usepackage{amsmath, amssymb, amsfonts}
\usepackage{graphicx}
\usepackage{adjustbox}

\hypersetup{hidelinks}

% Necesario para los .md convertidos con pandoc
\providecommand{\gt}{>}
\providecommand{\lt}{<}
\providecommand{\tightlist}{\setlength{\itemsep}{0pt}\setlength{\parskip}{0pt}}
\DeclareUnicodeCharacter{2264}{\ensuremath{\le}}
\DeclareUnicodeCharacter{2265}{\ensuremath{\ge}}
\DeclareUnicodeCharacter{2260}{\ensuremath{\ne}}
\DeclareUnicodeCharacter{2248}{\ensuremath{\approx}}
\DeclareUnicodeCharacter{2261}{\ensuremath{\equiv}}
\DeclareUnicodeCharacter{2192}{\ensuremath{\to}}
\DeclareUnicodeCharacter{2190}{\ensuremath{\leftarrow}}
\DeclareUnicodeCharacter{21D2}{\ensuremath{\Rightarrow}}
\DeclareUnicodeCharacter{00D7}{\ensuremath{\times}}
\DeclareUnicodeCharacter{00F7}{\ensuremath{\div}}
\DeclareUnicodeCharacter{22C5}{\ensuremath{\cdot}}
\DeclareUnicodeCharacter{2211}{\ensuremath{\sum}}
\DeclareUnicodeCharacter{220F}{\ensuremath{\prod}}
\DeclareUnicodeCharacter{221E}{\ensuremath{\infty}}
\DeclareUnicodeCharacter{2208}{\ensuremath{\in}}
\DeclareUnicodeCharacter{2209}{\ensuremath{\notin}}
\DeclareUnicodeCharacter{2229}{\ensuremath{\cap}}
\DeclareUnicodeCharacter{222A}{\ensuremath{\cup}}
\DeclareUnicodeCharacter{2212}{\ensuremath{-}}
\DeclareUnicodeCharacter{221A}{\ensuremath{\surd}}
\DeclareUnicodeCharacter{2026}{\ldots}

\renewcommand*{\ttdefault}{pcr}
\renewcommand\cftsecfont{\fontsize{10}{9}\bfseries}
\renewcommand\cftsecpagefont{\fontsize{10}{9}\mdseries}
\renewcommand\cftsubsecfont{\fontsize{10}{6}\mdseries}
\renewcommand\cftsubsecpagefont{\fontsize{10}{6}\mdseries}
\renewcommand\cftsecafterpnum{\vspace{-.5ex}}
\renewcommand\cftsubsecafterpnum{\vspace{0ex}}

\newcommand{\Section}[1]{
    \section*{#1}
    \addcontentsline{toc}{subsection}{#1}
    \vspace{-3mm}
}

\newcommand{\Topic}[1]{\textbf{#1}}

\lstdefinestyle{shared}{
    belowcaptionskip=1\baselineskip,
    breaklines=true,
    xleftmargin=0pt,
    showstringspaces=false,
    basicstyle=\fontsize{8}{5}\ttfamily,
    aboveskip=0pt,
    belowskip=0pt,
    tabsize=2,
    % Tildes y eñes en los comentarios del código
    literate=
      {á}{{\'a}}1 {é}{{\'e}}1 {í}{{\'i}}1 {ó}{{\'o}}1 {ú}{{\'u}}1
      {Á}{{\'A}}1 {É}{{\'E}}1 {Í}{{\'I}}1 {Ó}{{\'O}}1 {Ú}{{\'U}}1
      {ñ}{{\~n}}1 {Ñ}{{\~N}}1 {ü}{{\"u}}1 {Ü}{{\"U}}1
      {¿}{{?`}}1 {¡}{{!`}}1,
}
\lstdefinestyle{cpp}{
    style=shared,
    language=C++,
    keywordstyle=\bfseries\color{green!40!black},
    commentstyle=\itshape\color{blue!90!black},
    identifierstyle=\color{black},
    stringstyle=\color{red!90!black},
}
\lstdefinestyle{py}{
    style=shared,
    language=Python,
    keywordstyle=\bfseries\color{green!40!black},
    commentstyle=\itshape\color{purple!40!black},
    identifierstyle=\color{blue},
    stringstyle=\color{orange},
}
\lstdefinestyle{java}{
    style=shared,
    language=Java,
    keywordstyle=\bfseries\color{green!40!black},
    commentstyle=\itshape\color{purple!40!black},
    identifierstyle=\color{blue},
    stringstyle=\color{red},
}
\lstdefinestyle{txt}{
    style=shared,
}
\lstset{style=shared, escapechar=@}


\title{\vspace{-4ex}\Large{CodeDeGrace ICPC Team Notebook}}
\author{}
\date{}

% Configure headers
\fancyhf{}
\fancyhead[L]{CodeDeGrace - Universidad Tecnologica de Pereira}
\fancyhead[C]{%
  \ifthenelse{\equal{\firstleftmark}{\lastleftmark}}%
    {\leftmark}%
    {\firstleftmark\ \& \lastleftmark}%
}
\fancyhead[R]{\thepage}
\renewcommand{\headrulewidth}{0.4pt}

% Customize section marks to include section number
\renewcommand{\sectionmark}[1]{\markboth{\thesection. #1}{}}

\begin{document}
\pagestyle{fancy}
\setlength{\columnseprule}{0.4pt}
\setlength{\columnsep}{10pt}

% Pagina 1: solo titulo y contenidos
\begin{multicols}{3}
\maketitle
\vspace{-13ex}
\tableofcontents
\end{multicols}

% Salto de pagina: el template y el resto empiezan en la pagina 2
\clearpage

\setlength{\columnsep}{10pt}
\begin{multicols}{3}
\section{Template}
\subsection{Template}
\raggedbottom\lstinputlisting[style=cpp]{temp_template.h}
\hrulefill

\section{Data-structures}
\subsection{Segment tree}
\raggedbottom\lstinputlisting[style=cpp]{temp_Data-structures-SegTree.cpp}
\hrulefill
\subsection{Segment tree iterativo}
\raggedbottom\lstinputlisting[style=cpp]{temp_Data-structures-SegTree_iterativo.cpp}
\hrulefill
\subsection{Segment tree tricks}
\raggedbottom\lstinputlisting[style=cpp]{temp_Data-structures-SegTree_tricks.cpp}
\hrulefill
\subsection{Segment tree with structures}
\raggedbottom\lstinputlisting[style=cpp]{temp_Data-structures-SegTree_with_structures.cpp}
\hrulefill
\subsection{Segment tree segmento con mayor suma}
\raggedbottom\lstinputlisting[style=cpp]{temp_Data-structures-SegTree_segmento_con_mayor_suma.cpp}
\hrulefill
\subsection{Merge Sort Tree}
\raggedbottom\lstinputlisting[style=cpp]{temp_Data-structures-Merge_sort_tree.cpp}
\hrulefill
\subsection{Segment tree lazy}
\raggedbottom\lstinputlisting[style=cpp]{temp_Data-structures-SegTree_lazy.cpp}
\hrulefill
\subsection{Segment tree lazy some operations}
\raggedbottom\lstinputlisting[style=cpp]{temp_Data-structures-SegTree_lazy_some_operations.cpp}
\hrulefill
\subsection{Segment tree lazy arithmetic progression}
\raggedbottom\lstinputlisting[style=cpp]{temp_Data-structures-SegTree_lazy_arithmetic_progression.cpp}
\hrulefill
\subsection{Segment tree lazy segmento con mayor suma}
\raggedbottom\lstinputlisting[style=cpp]{temp_Data-structures-SegTree_lazy_segmento_con_mayor_suma.cpp}
\hrulefill
\subsection{Segment tree persistent}
\raggedbottom\lstinputlisting[style=cpp]{temp_Data-structures-SegTree_persistent.cpp}
\hrulefill
\subsection{Segment tree sparse}
\raggedbottom\lstinputlisting[style=cpp]{temp_Data-structures-SegTree_sparse.h}
\hrulefill
\subsection{Segment tree sparse lazy}
\raggedbottom\lstinputlisting[style=cpp]{temp_Data-structures-SegTree_sparse_lazy.cpp}
\hrulefill
\subsection{Segment tree 2D}
\raggedbottom\lstinputlisting[style=cpp]{temp_Data-structures-SegTree_2D.cpp}
\hrulefill
\subsection{Segment tree 2D sparse}
\raggedbottom\lstinputlisting[style=cpp]{temp_Data-structures-SegTree_2D_sparse.cpp}
\hrulefill
\subsection{Fenwick Tree}
\raggedbottom\lstinputlisting[style=cpp]{temp_Data-structures-FenwickTree.cpp}
\hrulefill
\subsection{Sparse Table}
\raggedbottom\lstinputlisting[style=cpp]{temp_Data-structures-SparseTable.cpp}
\hrulefill
\subsection{SQRT Decomposition}
\raggedbottom\lstinputlisting[style=cpp]{temp_Data-structures-SQRT_Decomposition.cpp}
\hrulefill
\subsection{Mo's Algorithm}
\raggedbottom\lstinputlisting[style=cpp]{temp_Data-structures-Mo_algorithm.cpp}
\hrulefill
\subsection{Rectangular partial sums}
\raggedbottom\lstinputlisting[style=cpp]{temp_Data-structures-Rectangular_partial_sums.cpp}
\hrulefill
\subsection{Kadane}
\raggedbottom\lstinputlisting[style=cpp]{temp_Data-structures-Kadane.cpp}
\hrulefill

\section{Graphs}
\subsection{Dijkstra}
\raggedbottom\lstinputlisting[style=cpp]{temp_Graphs-Dijkstra.cpp}
\hrulefill
\subsection{Bellman Ford}
\raggedbottom\lstinputlisting[style=cpp]{temp_Graphs-Bellman_Ford.cpp}
\hrulefill
\subsection{SPFA}
\raggedbottom\lstinputlisting[style=cpp]{temp_Graphs-SPFA.cpp}
\hrulefill
\subsection{Floyd Warshall}
\raggedbottom\lstinputlisting[style=cpp]{temp_Graphs-Floyd_Warshal.cpp}
\hrulefill
\subsection{BFS 0-1}
\raggedbottom\lstinputlisting[style=cpp]{temp_Graphs-BFS_01.cpp}
\hrulefill
\subsection{Find edges and vertexes in shortest paths}
\raggedbottom\lstinputlisting[style=cpp]{temp_Graphs-Find_edges_and_vertexes_in_shortest_path.cpp}
\hrulefill
\subsection{LCA BFS}
\raggedbottom\lstinputlisting[style=cpp]{temp_Graphs-LCA_BFS.cpp}
\hrulefill
\subsection{DSU}
\raggedbottom\lstinputlisting[style=cpp]{temp_Graphs-DSU.cpp}
\hrulefill
\subsection{DSU bipartition test}
\raggedbottom\lstinputlisting[style=cpp]{temp_Graphs-DSU_bipartition_test.cpp}
\hrulefill
\subsection{DSU rollback}
\raggedbottom\lstinputlisting[style=cpp]{temp_Graphs-DSU_rollback.cpp}
\hrulefill
\subsection{Dynamic connectivity offline}
\raggedbottom\lstinputlisting[style=cpp]{temp_Graphs-Dynamic_connectivity_offline.cpp}
\hrulefill
\subsection{Kruskal}
\raggedbottom\lstinputlisting[style=cpp]{temp_Graphs-Kruskal.cpp}
\hrulefill
\subsection{Prim}
\raggedbottom\lstinputlisting[style=cpp]{temp_Graphs-Prim.cpp}
\hrulefill
\subsection{Tarjan SCC}
\raggedbottom\lstinputlisting[style=cpp]{temp_Graphs-Tarjan_SCC.cpp}
\hrulefill
\subsection{Toposort}
\raggedbottom\lstinputlisting[style=cpp]{temp_Graphs-Toposort.cpp}
\hrulefill
\subsection{Bridges}
\raggedbottom\lstinputlisting[style=cpp]{temp_Graphs-Bridges.cpp}
\hrulefill
\subsection{Articulation points}
\raggedbottom\lstinputlisting[style=cpp]{temp_Graphs-Articulation_points.cpp}
\hrulefill
\subsection{TSP}
\raggedbottom\lstinputlisting[style=cpp]{temp_Graphs-TSP.cpp}
\hrulefill
\subsection{Counting cycles}
\raggedbottom\lstinputlisting[style=cpp]{temp_Graphs-Counting_cycles.cpp}
\hrulefill
\subsection{Counting paths}
\raggedbottom\lstinputlisting[style=cpp]{temp_Graphs-Counting_paths.cpp}
\hrulefill
\subsection{Eulerian path}
\raggedbottom\lstinputlisting[style=cpp]{temp_Graphs-Eulerian_path.cpp}
\hrulefill
\subsection{Grafo funcional Floyd}
\raggedbottom\lstinputlisting[style=cpp]{temp_Graphs-Grafo_funcional_Floyd.cpp}
\hrulefill
\subsection{Grafo funcional count steps}
\raggedbottom\lstinputlisting[style=cpp]{temp_Graphs-Grafo_funcional_count_steps.cpp}
\hrulefill
\subsection{Grafo funcional kth step}
\raggedbottom\lstinputlisting[style=cpp]{temp_Graphs-Grafo_funcional_kth_step.cpp}
\hrulefill
\subsection{Find faces grafo plano}
\raggedbottom\lstinputlisting[style=cpp]{temp_Graphs-Find_faces_grafo_plano.cpp}
\hrulefill

\section{Trees}
\subsection{Diameter}
\raggedbottom\lstinputlisting[style=cpp]{temp_Trees-Diameter.cpp}
\hrulefill
\subsection{Euler Tour}
\raggedbottom\lstinputlisting[style=cpp]{temp_Trees-Euler_Tour.cpp}
\hrulefill
\subsection{LCA Euler Tour}
\raggedbottom\lstinputlisting[style=cpp]{temp_Trees-LCA_Euler_Tour.cpp}
\hrulefill
\subsection{LCA and Binary Lifting}
\raggedbottom\lstinputlisting[style=cpp]{temp_Trees-LCA_and_BinaryLifting.cpp}
\hrulefill
\subsection{Heavy light decomposition}
\raggedbottom\lstinputlisting[style=cpp]{temp_Trees-Heavy_light_decomposition.cpp}
\hrulefill
\subsection{Max clique}
\raggedbottom\lstinputlisting[style=cpp]{temp_Trees-Max_clique.cpp}
\hrulefill
\subsection{Minimum vertex cover}
\raggedbottom\lstinputlisting[style=cpp]{temp_Trees-Minimum_vertex_cover.cpp}
\hrulefill

\section{Flow}
\subsection{Dinic}
\raggedbottom\lstinputlisting[style=cpp]{temp_Flow-MaxFlow_Dinic.cpp}
\hrulefill
\subsection{Push Relabel}
\raggedbottom\lstinputlisting[style=cpp]{temp_Flow-MaxFlow_PushRelabel.cpp}
\hrulefill
\subsection{Hungarian}
\raggedbottom\lstinputlisting[style=cpp]{temp_Flow-Hungarian.cpp}
\hrulefill
\subsection{Min Cost Max Flow}
\raggedbottom\lstinputlisting[style=cpp]{temp_Flow-MinCostMaxFlow.cpp}
\hrulefill
\subsection{Blossom}
\raggedbottom\lstinputlisting[style=cpp]{temp_Flow-Blossom.cpp}
\hrulefill

\section{DP}
\subsection{Divide and Conquer}
\raggedbottom\lstinputlisting[style=cpp]{temp_DP-D-Q.cpp}
\hrulefill
\subsection{Knuth}
\raggedbottom\lstinputlisting[style=cpp]{temp_DP-Knuth.cpp}
\hrulefill
\subsection{Convex Hull Trick}
\raggedbottom\lstinputlisting[style=cpp]{temp_DP-CHT.cpp}
\hrulefill
\subsection{Convex Hull Trick dinamico}
\raggedbottom\lstinputlisting[style=cpp]{temp_DP-CHT_dynamic.cpp}
\hrulefill
\subsection{Digit DP}
\raggedbottom\lstinputlisting[style=cpp]{temp_DP-Digit_DP.cpp}
\hrulefill
\subsection{LIS}
\raggedbottom\lstinputlisting[style=cpp]{temp_DP-LIS.cpp}
\hrulefill
\subsection{Recorrer submasks}
\raggedbottom\lstinputlisting[style=cpp]{temp_DP-Recorrer_submasks.cpp}
\hrulefill

\section{Math}
\subsection{Combinatory}
\begingroup\setlength{\parindent}{0pt}\setlength{\parskip}{2pt}\input{temp_Math-Combinatory.tex}\par\endgroup
\hrulefill
\subsection{Coeficientes binomiales}
\raggedbottom\lstinputlisting[style=cpp]{temp_Math-Coeficientes_binomiales.cpp}
\hrulefill
\subsection{Power of factorial divisor}
\raggedbottom\lstinputlisting[style=cpp]{temp_Math-Power_of_factorial_divisor.cpp}
\hrulefill
\subsection{Factorial MOD}
\raggedbottom\lstinputlisting[style=cpp]{temp_Math-FactorialMOD.cpp}
\hrulefill
\subsection{Matriz}
\raggedbottom\lstinputlisting[style=cpp]{temp_Math-Matriz.cpp}
\hrulefill
\subsection{Probability}
\begingroup\setlength{\parindent}{0pt}\setlength{\parskip}{2pt}\input{temp_Math-Probability.tex}\par\endgroup
\hrulefill
\subsection{FFT}
\raggedbottom\lstinputlisting[style=cpp]{temp_Math-FFT.cpp}
\hrulefill
\subsection{Random}
\raggedbottom\lstinputlisting[style=cpp]{temp_Math-Random.cpp}
\hrulefill
\subsection{Composicion de funciones}
\raggedbottom\lstinputlisting[style=cpp]{temp_Math-Composicion_funciones.cpp}
\hrulefill
\subsection{Numeros practicos}
\raggedbottom\lstinputlisting[style=cpp]{temp_Math-Numeros-Practicos.cpp}
\hrulefill

\section{Number Theory}
\subsection{Binpow}
\raggedbottom\lstinputlisting[style=cpp]{temp_Number-Theory-Binpow.cpp}
\hrulefill
\subsection{Inverso modular}
\raggedbottom\lstinputlisting[style=cpp]{temp_Number-Theory-Inverso_modular.cpp}
\hrulefill
\subsection{Descomposicion prima}
\raggedbottom\lstinputlisting[style=cpp]{temp_Number-Theory-Descomposicion_prima.cpp}
\hrulefill
\subsection{Sieve of Eratosthenes}
\raggedbottom\lstinputlisting[style=cpp]{temp_Number-Theory-Sieve_erastosthenes.cpp}
\hrulefill
\subsection{Segmented sieve}
\raggedbottom\lstinputlisting[style=cpp]{temp_Number-Theory-Segmented_sieve.cpp}
\hrulefill
\subsection{Divisores de cada numero}
\raggedbottom\lstinputlisting[style=cpp]{temp_Number-Theory-Divisores_de_cada_numero.cpp}
\hrulefill
\subsection{Miller Rabin}
\raggedbottom\lstinputlisting[style=cpp]{temp_Number-Theory-Miller_rabin.cpp}
\hrulefill
\subsection{Pollard rho}
\raggedbottom\lstinputlisting[style=cpp]{temp_Number-Theory-Pollard_rho.cpp}
\hrulefill
\subsection{Funcion de Mobius}
\raggedbottom\lstinputlisting[style=cpp]{temp_Number-Theory-Funcion_mobius.cpp}
\hrulefill
\subsection{Funcion de Euler}
\raggedbottom\lstinputlisting[style=cpp]{temp_Number-Theory-Funcion_euler.cpp}
\hrulefill
\subsection{Mulpow}
\raggedbottom\lstinputlisting[style=cpp]{temp_Number-Theory-Mulpow.cpp}
\hrulefill
\subsection{Euclides}
\raggedbottom\lstinputlisting[style=cpp]{temp_Number-Theory-Euclides.cpp}
\hrulefill
\subsection{Diofanticas}
\raggedbottom\lstinputlisting[style=cpp]{temp_Number-Theory-Diofanticas.cpp}
\hrulefill
\subsection{Chinese remainder theorem}
\raggedbottom\lstinputlisting[style=cpp]{temp_Number-Theory-Chinese_remainder_theorem.cpp}
\hrulefill
\subsection{Fibonacci code}
\raggedbottom\lstinputlisting[style=cpp]{temp_Number-Theory-Fibonacci_code.cpp}
\hrulefill
\subsection{Teoria}
\begingroup\setlength{\parindent}{0pt}\setlength{\parskip}{2pt}\input{temp_Number-Theory-Teoria.tex}\par\endgroup
\hrulefill

\section{Miscelanea}
\subsection{Compresion}
\raggedbottom\lstinputlisting[style=cpp]{temp_Miscelanea-Compresion.cpp}
\hrulefill
\subsection{Biseccion}
\raggedbottom\lstinputlisting[style=cpp]{temp_Miscelanea-Biseccion.cpp}
\hrulefill
\subsection{Busqueda ternaria}
\raggedbottom\lstinputlisting[style=cpp]{temp_Miscelanea-Busquedad_ternaria.cpp}
\hrulefill
\subsection{Bits}
\raggedbottom\lstinputlisting[style=cpp]{temp_Miscelanea-Bits.cpp}
\hrulefill
\subsection{Sliding window}
\raggedbottom\lstinputlisting[style=cpp]{temp_Miscelanea-Sliding_window.cpp}
\hrulefill
\subsection{Game theory}
\begingroup\setlength{\parindent}{0pt}\setlength{\parskip}{2pt}\input{temp_Miscelanea-Game_theory.tex}\par\endgroup
\hrulefill
\subsection{Games on arbitrary graphs}
\raggedbottom\lstinputlisting[style=cpp]{temp_Miscelanea-Games_on_arbitrary_graphs.cpp}
\hrulefill
\subsection{Arbol cartesiano}
\raggedbottom\lstinputlisting[style=cpp]{temp_Miscelanea-Arbol_cartesiano.cpp}
\hrulefill
\subsection{Combinaciones longitud k}
\raggedbottom\lstinputlisting[style=cpp]{temp_Miscelanea-Combinaciones_longitud_k.cpp}
\hrulefill
\subsection{Josefo}
\raggedbottom\lstinputlisting[style=cpp]{temp_Miscelanea-Josefo.cpp}
\hrulefill
\subsection{Hashing XOR}
\raggedbottom\lstinputlisting[style=cpp]{temp_Miscelanea-Hashing_XOR.cpp}
\hrulefill
\subsection{Python}
\raggedbottom\lstinputlisting[style=py]{temp_Miscelanea-Python.py}
\hrulefill

\section{C++}
\subsection{Bitset}
\raggedbottom\lstinputlisting[style=cpp]{temp_C---Bitset.cpp}
\hrulefill
\subsection{Enteros de 128 bits}
\raggedbottom\lstinputlisting[style=cpp]{temp_C---Enteros_128.cpp}
\hrulefill
\subsection{Imprimir decimales}
\raggedbottom\lstinputlisting[style=cpp]{temp_C---Imprimir_decimales.cpp}
\hrulefill
\subsection{Deque}
\raggedbottom\lstinputlisting[style=cpp]{temp_C---Deque.cpp}
\hrulefill
\subsection{Indexed set}
\raggedbottom\lstinputlisting[style=cpp]{temp_C---Indexed_set.cpp}
\hrulefill
\subsection{Listas}
\raggedbottom\lstinputlisting[style=cpp]{temp_C---Listas.cpp}
\hrulefill
\subsection{Map}
\raggedbottom\lstinputlisting[style=cpp]{temp_C---Map.cpp}
\hrulefill
\subsection{Metodos de ordenamiento}
\raggedbottom\lstinputlisting[style=cpp]{temp_C---Metodos_ordenar.cpp}
\hrulefill
\subsection{Permutaciones}
\raggedbottom\lstinputlisting[style=cpp]{temp_C---Permutaciones.cpp}
\hrulefill

\section{Strings}
\subsection{Hashing}
\raggedbottom\lstinputlisting[style=cpp]{temp_Strings-Hashing.cpp}
\hrulefill
\subsection{Minimal Rotation}
\raggedbottom\lstinputlisting[style=cpp]{temp_Strings-Minimal_rotation.cpp}
\hrulefill
\subsection{String Matching}
\raggedbottom\lstinputlisting[style=cpp]{temp_Strings-String_matching.cpp}
\hrulefill
\subsection{Trie}
\raggedbottom\lstinputlisting[style=cpp]{temp_Strings-Trie.cpp}
\hrulefill

\section{Geometria}
\subsection{Geo Struct}
\raggedbottom\lstinputlisting[style=cpp]{temp_Geometria-GeoStruct.cpp}
\hrulefill
\subsection{Convex Hull}
\raggedbottom\lstinputlisting[style=cpp]{temp_Geometria-ConvexHull.cpp}
\hrulefill
\subsection{Intersection Points}
\raggedbottom\lstinputlisting[style=cpp]{temp_Geometria-IntersectionPoints.cpp}
\hrulefill
\subsection{Line Segment Intersection}
\raggedbottom\lstinputlisting[style=cpp]{temp_Geometria-LineSegmentIntersection.cpp}
\hrulefill
\subsection{Segment Intersection}
\raggedbottom\lstinputlisting[style=cpp]{temp_Geometria-SegmentIntersection.cpp}
\hrulefill
\subsection{Point in Polygon}
\raggedbottom\lstinputlisting[style=cpp]{temp_Geometria-PointinPolygon.cpp}
\hrulefill
\subsection{Polygons Area}
\raggedbottom\lstinputlisting[style=cpp]{temp_Geometria-PolygonsArea.cpp}
\hrulefill
\subsection{Li Chao Tree}
\raggedbottom\lstinputlisting[style=cpp]{temp_Geometria-LiChaoTree.cpp}
\hrulefill
\subsection{Gauss (UNTESTED)}
\raggedbottom\lstinputlisting[style=cpp]{temp_Geometria-Gauss.cpp}
\hrulefill


\end{multicols}

% Solo si tienes un Theory.tex aparte (tus .md ya van dentro de contents.tex):
% \setlength{\columnseprule}{0pt}
% \setlength{\columnsep}{10pt}
% \begin{multicols}{2}
% {\fontsize{8}{6}\selectfont
% \input{Theory.tex}
% }
% \end{multicols}

\end{document}